\documentclass[12pt,a4paper]{article}
% !TeX program = pdflatex
% !BIB program = bibtex
\usepackage{graphicx}
\usepackage{float}
\usepackage{amsmath}
\usepackage{booktabs}
\usepackage{geometry}
\usepackage[round]{natbib}
\usepackage{xcolor}
\usepackage{setspace}
\usepackage{tocloft}
\usepackage{titlesec}
\usepackage{rotating}
\usepackage[none]{hyphenat}
\usepackage[skip=6pt]{caption}
\usepackage{tocloft}
\setlength{\cftsubsecnumwidth}{3em}
\usepackage{listings}
\usepackage{xcolor}

\lstset{
  basicstyle=\ttfamily\small,
  breaklines=true,
  frame=single
}

\geometry{margin=1in}

\setcounter{secnumdepth}{3}
\setcounter{tocdepth}{3}
\setlength{\cftbeforesecskip}{5pt}
\setlength{\cftbeforesubsecskip}{2pt}
\newcommand{\sectionbreak}{\clearpage}

\begin{document}
% =============================================================
\begin{titlepage}
\centering

\vspace*{\fill}   % pushes content to vertical center

{\Large \textbf{Analyzing Global Music Trends Using Spotify Audio Features:}\\[0.4cm]
\textbf{A Data-Driven Analysis in R}}

\vspace{1.2cm}

{\large Shristi}\\[0.2cm]
\textit{Author}

\vspace{0.4cm}

B.Tech, Computer Science and Engineering with Specialization in Data Science\\[0.3cm]
Indian Institute of Information Technology Una\\
Himachal Pradesh, India

\vspace{0.8cm}
{\large Dr. Prince Sharma}\\[0.2cm]
\textit{Co-Author}

\vspace*{\fill}   % balances bottom

\end{titlepage}
\setcounter{page}{1}
% =============================================================

\newpage
\section*{Abstract}

\begin{spacing}{1.2}
\noindent
The widespread adoption of digital music streaming platforms has created unprecedented opportunities to analyze large-scale music consumption data. This study investigates global music trends by analyzing Spotify track data using statistical and visualization techniques implemented in R. The objective is to explore how various audio features, including danceability, energy, valence, acousticness, and tempo, relate to song popularity and how these characteristics have evolved across different release periods.

\vspace{0.5em}

The analysis uses a publicly available dataset of Spotify tracks and follows a structured data analytics workflow comprising data preprocessing, exploratory data analysis, correlation analysis, and feature-based trend visualization. Using R packages such as the tidyverse, ggplot2, and dplyr, the study examines relationships between musical attributes and listener engagement and identifies patterns that characterise highly popular tracks.

\vspace{0.5em}

Visual analytics techniques are employed to reveal temporal changes in audio feature distributions and highlight potential associations between musical structure and audience preferences. Findings show that danceability is the strongest positive predictor, 
while speechiness, valence, and instrumentalness show significant 
negative associations with popularity.

\vspace{0.5em}

This study demonstrates the effectiveness of R as a flexible platform for large-scale music data analysis and illustrates how data mining and visualization methods can provide valuable insights into modern music consumption patterns. The findings contribute to a broader understanding of how data-driven approaches can be applied to analyze trends in the digital music ecosystem.
\end{spacing}

\clearpage

% ===== CONTENTS =====
\clearpage
\begingroup
\setstretch{1.1}
\tableofcontents
\setlength{\cftbeforesecskip}{6pt}
\endgroup
\clearpage
% =============================================================
% 
%=============================================================
\section{Introduction}
%=============================================================

\subsection{Background of Music Streaming Platforms}

The global music industry has undergone a significant transformation over
the past decade. Digital streaming platforms have largely replaced
physical sales and downloads as the primary mode of music consumption
worldwide. Among these platforms, Spotify, launched in 2008, is one of the leading services, offering access to millions of tracks
across a wide range of genres and languages \citep{spotify2023}. Recent
estimates suggest that the platform hosts over 100 million tracks and
serves hundreds of millions of active users globally \citep{ifpi2023}.

Unlike traditional broadcast media, streaming platforms generate
continuous and highly detailed data on user behaviour. Every
interaction---such as plays, skips, saves, and playlist additions—is
recorded and stored. Spotify provides access to a subset of this data
through its Web API, which includes several computed audio features.
These features, such as \textit{danceability}, \textit{energy}, and
\textit{valence}, represent different musical characteristics in
numerical form. Additionally, a \textit{popularity} score, ranging from
0 to 100, indicates the relative level of listener engagement for each
track \citep{spotifyapi2023}.

\subsection{Importance of Data-Driven Music Analysis}

The availability of large-scale audio feature data enables a more
systematic and data-driven approach to understanding music trends.
Traditional music analysis methods are often qualitative and limited to
small datasets, making it difficult to generalize findings. In contrast,
modern data analysis techniques allow researchers to identify patterns
across hundreds of thousands of tracks simultaneously \citep{schedl2014}.

Understanding the relationship between audio features and popularity has
important practical applications. Artists and producers can use these
insights to create music that better aligns with audience preferences,
while streaming platforms can improve their recommendation systems.
From a research perspective, this approach contributes to the field of
computational musicology by examining whether measurable audio features
have a statistically significant impact on song popularity, beyond
external factors such as marketing, timing, or social influence
\citep{lamere2014,interiano2018}.

In this study, a hierarchical regression model is applied to a large
Spotify dataset obtained from Kaggle \citep{kaggle_spotify}. The analysis is carried out using
the R programming language \citep{rcoreteam2023}, ensuring a structured, transparent, and
reproducible workflow. The study aims to identify key audio features that
influence song popularity and to understand how these relationships can
be interpreted in a real-world context.
% =============================================================
% =============================================================
\section{Problem Statement}
% =============================================================

Although Spotify audio features are widely used in exploratory data
analysis, there is limited work that systematically evaluates their
combined and incremental impact on song popularity using structured
statistical models. Most existing studies rely on simple correlation
analysis, single-model regression, or machine learning approaches that
focus primarily on prediction accuracy rather than interpretability.

As a result, there is a lack of clear understanding regarding how
different groups of audio features contribute to popularity when
considered together in a step-by-step manner.

To address this gap, the present study focuses on the following research
question:

\begin{quote}
\textbf{Which Spotify audio features have a statistically significant
impact on song popularity, and how does the inclusion of different
feature groups improve the explanatory power of the model?}
\end{quote}

By applying a hierarchical regression approach, this study aims to
analyze the individual as well as combined effects of multiple audio
features, providing a clearer and more interpretable understanding of
their relationship with song popularity.
% =============================================================
\section{Objectives of the Study}
% =============================================================

The primary objective of this study is to analyze the relationship
between Spotify audio features and song popularity using a structured
data-driven approach.

The specific objectives are as follows:

\begin{enumerate}
  \item To load, clean, and prepare the Spotify tracks dataset using
        reproducible techniques in R.
        
  \item To perform exploratory data analysis (EDA) to understand
        the distribution and behaviour of key audio features and the
        popularity variable.
        
  \item To examine the relationships between different audio features
        and song popularity using correlation analysis and visual
        exploration.
        
  \item To develop a hierarchical regression model by introducing
        features in multiple blocks and evaluating their incremental
        contribution to explaining song popularity.
        
  \item To assess model performance using statistical measures such as
        $R^2$, Adjusted $R^2$, and AIC.
        
  \item To interpret regression coefficients along with confidence
        intervals to understand the direction and strength of
        feature impact.
        
  \item To derive meaningful insights and conclusions based on the
        statistical results and visual analysis.
\end{enumerate}
% =============================================================
\section{Dataset Description}
% =============================================================

\subsection{Data Source}

The dataset used in this study is \texttt{spotify\_tracks.csv}, which is
publicly available on Kaggle. It has been compiled using the Spotify Web
API and contains one record per track, along with multiple audio feature
descriptors and metadata fields \citep{kaggle_spotify}.

\subsection{Dataset Overview}

After loading and applying the data cleaning process described in
Section~\ref{sec:preprocessing}, the final dataset consisted of
\textbf{114{,}000 observations} and \textbf{22 variables}. No missing
values were found in the nine key variables used for analysis, ensuring a
complete and reliable dataset for regression modelling.

\subsection{Description of Key Features}

The main variables used in the analysis are summarized in
Table 1. These features represent different musical
characteristics of a track and are provided directly by the Spotify API.

\begin{table}[ht]
\centering

\begin{tabular}{@{}lll@{}}
\toprule
\textbf{Variable}       & \textbf{Range} & \textbf{Description} \\
\midrule
\texttt{danceability}   & 0 -- 1 & Suitability of a track for dancing \\
\texttt{energy}         & 0 -- 1 & Perceived intensity and activity level \\
\texttt{valence}        & 0 -- 1 & Musical positiveness or mood \\
\texttt{acousticness}   & 0 -- 1 & Likelihood that the track is acoustic \\
\texttt{instrumentalness} & 0 -- 1 & Likelihood of absence of vocals \\
\texttt{speechiness}    & 0 -- 1 & Presence of spoken words in the track \\
\texttt{tempo}          & BPM    & Estimated tempo of the track \\
\texttt{loudness}       & dBFS   & Overall loudness (typically negative) \\
\midrule
\texttt{popularity}     & 0 -- 100 & \textbf{Target (dependent) variable} \\
\bottomrule
\end{tabular}

\vspace{0.3cm}

{\centering \textit{Table 1: Description of key Spotify audio features used in the analysis}\par}

\label{tab:features}

\end{table}

\subsection{Target Variable: Popularity}

The \texttt{popularity} score is a Spotify-generated integer ranging from
0 to 100, where higher values indicate greater listener engagement and
recent streaming activity \citep{spotifyapi2023}. 

In the cleaned dataset, the popularity distribution is slightly
right-skewed, with a mean of \textbf{33.24} and a standard deviation of
\textbf{22.31}. The median value is 35, with observations ranging from 0
to 100. The interquartile range lies between 17 and 50, indicating a wide
variation in popularity levels and a relatively large number of
low-popularity tracks.
% =============================================================
\section{Methodology}
% =============================================================

\subsection{Data Collection}

The dataset used in this study was obtained from Kaggle in the form of a
CSV file (\texttt{spotify\_tracks.csv}). It was originally compiled using
the Spotify Web API and contains detailed information about individual
tracks along with their corresponding audio features \citep{kaggle_spotify}.

\subsection{Data Loading in R}
The dataset was loaded into the R environment using functions from the
\texttt{tidyverse} package. The \texttt{read\_csv()} function was used to
import the data efficiently while preserving its structure.

\begin{verbatim}
library(tidyverse)
library(janitor)

df_raw <- read_csv("spotify_tracks.csv",
                   show_col_types = FALSE)
\end{verbatim}

\subsection{Data Preprocessing}
\label{sec:preprocessing}

Data preprocessing was performed to ensure accuracy and consistency of
the dataset. Column names were standardised into \texttt{snake\_case}
format using the \texttt{janitor} package \citep{rcoreteam2023}. Missing values in key variables
were removed using appropriate filtering techniques, and duplicate rows
were eliminated to avoid redundancy in the analysis.

\subsection{Exploratory Data Analysis (EDA)}

Exploratory Data Analysis (EDA) was conducted to understand the
distribution and relationships among variables. This included the
calculation of descriptive statistics such as mean, median, and
quartiles. Visual analysis was performed using histograms, scatter plots,
and correlation matrices to identify patterns and trends in the data.

\subsection{Statistical Analysis}

To examine the relationship between audio features and song popularity,
three hierarchical Ordinary Least Squares (OLS) regression models were
constructed: 
Specifically, Block 1 introduced rhythmic and mood features
(\texttt{danceability}, \texttt{energy}, \texttt{valence}),
Block 2 added acoustic texture features (\texttt{acousticness},
\texttt{instrumentalness}, \texttt{speechiness}), and Block 3
incorporated physical signal properties (\texttt{tempo}, \texttt{loudness}). Each model introduced additional groups of variables to
evaluate their incremental contribution to the explanation of popularity.

Model performance was assessed using statistical measures such as
$R^2$, Adjusted $R^2$, and Akaike Information Criterion (AIC).
Additionally, ANOVA tests were used to compare models and determine
whether the inclusion of new variables significantly improved the model.

\subsection{Data Visualization Approach}

Data visualisation was carried out using the \texttt{ggplot2} package in
R. Various plots such as histograms, scatter plots, and heatmaps were
created to effectively represent patterns in the data. These visualisations were combined into a single dashboard using the
\texttt{patchwork} package. A correlation heatmap was also generated
using the \texttt{corrplot} package to highlight relationships among
key variables.

% =============================================================
\section{Tools and Technologies Used}
% =============================================================

\subsection{Programming Language: R}

All analyses were carried out using R (version 4.4.1), an open-source
statistical programming language widely used for data analysis and
visualization in academic and industry settings \citep{rcoreteam2023}.

\subsection{Libraries and Packages}

The following R packages were used for data processing, analysis, and
visualization:

\begin{table}[ht]
\centering

\begin{tabular}{@{}lll@{}}
\toprule
\textbf{Package}   & \textbf{Version} & \textbf{Purpose} \\
\midrule
\texttt{tidyverse} & $\geq$ 2.0   & Data manipulation and visualization (\texttt{ggplot2}) \\
\texttt{janitor}   & $\geq$ 2.2   & Cleaning and standardizing column names \\
\texttt{broom}     & $\geq$ 1.0   & Converting model outputs into tidy format \\
\texttt{corrplot}  & $\geq$ 0.92  & Creating correlation heatmaps \\
\texttt{patchwork} & $\geq$ 1.2   & Combining multiple plots into one layout \\
\texttt{scales}    & $\geq$ 1.3   & Formatting axes and labels in plots \\
\texttt{viridis}   & $\geq$ 0.6   & Providing color-blind-friendly palettes \\
\bottomrule
\end{tabular}

\vspace{0.3cm}

{\centering \textit{Table 2: R packages used in this study}\par}

\label{tab:packages}

\end{table}
% =============================================================
\section{Data Preprocessing}
\label{sec:preprocessing}
% =============================================================

\subsection{Handling Missing Values}

After standardizing column names using \texttt{clean\_names()} from the \texttt{janitor} package
\citep{rcoreteam2023}, the
dataset was examined for missing values in the nine key variables used
for analysis. The inspection confirmed that there were \textbf{no missing
values} in either the predictor or the outcome variables.

To ensure data consistency, rows containing missing values (if any) were
removed using appropriate filtering techniques, resulting in a
complete-case dataset suitable for regression analysis.
\begin{verbatim}
df <- df |> drop_na(all_of(key_cols))
df <- df |> distinct()
\end{verbatim}
\subsection{Removing Duplicates}

Duplicate entries in the dataset were identified and removed using
standard data cleaning methods. After this step, the dataset retained
\textbf{114{,}000 unique track observations}, ensuring that each record
represented a distinct track.

\subsection{Data Type Verification}

The data types of all variables were examined to ensure correctness.
All audio feature variables were found to be numeric, while the
\texttt{popularity} variable was stored as an integer. Since the data
types were already appropriate, no further conversion was required.

\subsection{Feature Selection}

A total of nine variables were selected for the regression analysis
based on their theoretical relevance and prior research. The variable
\texttt{popularity} was used as the dependent (target) variable, while
the following features were used as predictors:

\texttt{danceability}, \texttt{energy}, \texttt{valence},
\texttt{acousticness}, \texttt{instrumentalness},
\texttt{speechiness}, \texttt{tempo}, and \texttt{loudness}.
% =============================================================
\section{Exploratory Data Analysis}
% =============================================================

\subsection{Summary Statistics}

Descriptive statistics for the nine key variables are presented in
Table 3.

\begin{table}[ht]
\centering

\begin{tabular}{@{}lrrrrr@{}}
\toprule
\textbf{Variable}   & \textbf{Min}  & \textbf{Q1}  & \textbf{Median}
                    & \textbf{Mean} & \textbf{Max} \\
\midrule
popularity          & 0.00  & 17.00  & 35.00  & 33.24  & 100.00 \\
danceability        & 0.000 & 0.456  & 0.580  & 0.567  & 0.985  \\
energy              & 0.000 & 0.472  & 0.685  & 0.641  & 1.000  \\
valence             & 0.000 & 0.260  & 0.464  & 0.474  & 0.995  \\
acousticness        & 0.000 & 0.017  & 0.169  & 0.315  & 0.996  \\
instrumentalness    & 0.000 & 0.000  & 0.000  & 0.156  & 1.000  \\
speechiness         & 0.000 & 0.036  & 0.049  & 0.085  & 0.965  \\
tempo (BPM)         & 0.00  & 99.22  & 122.02 & 122.15 & 243.37 \\
loudness (dBFS)     & $-$49.53 & $-$10.01 & $-$7.00 & $-$8.26 & 4.53 \\
\bottomrule
\end{tabular}

\vspace{0.3cm}

{\centering \textit{Table 3: Summary statistics for key variables ($n = 114{,}000$)}\par}

\label{tab:summary}

\end{table}

\subsection{Distribution of Popularity}

The popularity values range from 0 to 100, with a mean of 33.24 and a
standard deviation of 22.31. The median value is 35, indicating that most
tracks fall in the low-to-moderate popularity range.

A noticeable concentration of tracks at zero popularity suggests the
presence of inactive or less-streamed tracks. The histogram (Figure 1)
visually confirms this pattern, showing a slightly skewed distribution
with a large number of low-popularity observations.

\subsection{Distribution of Audio Features}

The analysis of individual audio features reveals the following patterns:

\begin{itemize}
  \item \textbf{Danceability} (mean = 0.567) shows moderate-to-high values,
        indicating that many tracks are rhythmically suitable for dancing.
        
  \item \textbf{Energy} (mean = 0.641) is relatively high, reflecting the
        dominance of energetic and amplified music in modern streaming.
        
  \item \textbf{Valence} (mean = 0.474) is close to the midpoint,
        suggesting a balanced mix of positive and neutral emotional tones.
        
  \item \textbf{Instrumentalness} (mean = 0.156) is highly skewed towards zero,
        indicating that most tracks contain vocals.
        
  \item \textbf{Loudness} ranges from $-49.53$ to $+4.53$ dBFS, with a mean
        of $-8.26$ dBFS, consistent with contemporary loudness standards.
\end{itemize}

\subsection{Genre-Based Analysis}

Although the regression models in this study focus on track-level
features, the dataset also includes a \texttt{track\_genre} variable.
Analyzing genre-wise differences in popularity can provide additional
insights and is identified as a potential extension of this work, as
discussed in Section~\ref{sec:future}.
% =============================================================
% =============================================================
\section{Statistical Analysis}
% =============================================================

\subsection{Correlation Analysis}

Pairwise Pearson correlations were computed among all nine variables.
The complete correlation matrix is presented in Table 4.

\begin{table}[H]
\centering
\small
\renewcommand{\arraystretch}{1.2}

\begin{tabular}{lrrrrrrrrr}
\toprule
 & pop. & dance. & energy & valence & acous & instr & speech & tempo & loud \\
\midrule
popularity       & 1.000 & 0.035 & 0.001 & -0.041 & -0.025 & -0.095 & -0.045 & 0.013 & 0.050 \\
danceability     & 0.035 & 1.000 & 0.134 & 0.477 & -0.172 & -0.186 & 0.109 & -0.050 & 0.259 \\
energy           & 0.001 & 0.134 & 1.000 & 0.259 & -0.734 & -0.182 & 0.143 & 0.248 & 0.762 \\
valence          & -0.041 & 0.477 & 0.259 & 1.000 & -0.107 & -0.324 & 0.037 & 0.078 & 0.280 \\
acousticness     & -0.025 & -0.172 & -0.734 & -0.107 & 1.000 & 0.104 & -0.002 & -0.208 & -0.590 \\
instrumentalness & -0.095 & -0.186 & -0.182 & -0.324 & 0.104 & 1.000 & -0.090 & -0.050 & -0.433 \\
speechiness      & -0.045 & 0.109 & 0.143 & 0.037 & -0.002 & -0.090 & 1.000 & 0.017 & 0.061 \\
tempo            & 0.013 & -0.050 & 0.248 & 0.078 & -0.208 & -0.050 & 0.017 & 1.000 & 0.212 \\
loudness         & 0.050 & 0.259 & 0.762 & 0.280 & -0.590 & -0.433 & 0.061 & 0.212 & 1.000 \\
\bottomrule
\end{tabular}

\vspace{0.4em}
\begin{center}
\textit{Table 4: Pearson correlation matrix for key variables}
\end{center}

\end{table}

\subsection{Correlation Matrix Interpretation}

The correlation values between individual audio features and popularity
are relatively small, indicating weak linear relationships. The strongest
correlation (in absolute terms) is observed between
\texttt{instrumentalness} and popularity ($r = -0.095$), followed by
\texttt{speechiness} ($r = -0.045$), \texttt{valence} ($r = -0.041$),
and \texttt{loudness} ($r = 0.050$).

Interestingly, \texttt{energy} shows almost no direct correlation with
popularity ($r = 0.001$). However, it appears as a statistically
significant predictor in the regression model. This indicates that the
effect of energy becomes more apparent when other variables are taken
into account.

Strong relationships are observed among predictor variables themselves.
For example, \texttt{energy} and \texttt{loudness} show a high positive
correlation ($r = 0.762$), while \texttt{energy} and
\texttt{acousticness} show a strong negative correlation
($r = -0.734$). These patterns suggest that some features capture similar
underlying musical properties.

\subsection{Feature Relationships}

Scatter plots of \texttt{danceability} versus popularity and
\texttt{energy} versus popularity (Figure~\ref{fig:dance_pop} and 
Figure~\ref{fig:energy_pop}) show weak but
noticeable trends. While higher danceability is slightly associated with
increased popularity, energy shows a less consistent relationship. In
both cases, a high degree of dispersion is observed, indicating that
audio features alone do not strongly determine popularity.

% =============================================================
% =============================================================
\section{Data Visualization}
% =============================================================

\subsection{Overview of Visual Analysis}

To better understand the relationships between audio features and song
popularity, multiple visualizations were created using \texttt{ggplot2}.
These plots help in identifying patterns, trends, and differences that
may not be immediately visible through numerical analysis alone.

% -------------------------------------------------------------

\subsection{Distribution of Song Popularity}

The distribution of the popularity variable was analysed using a
histogram with 50 bins. The popularity values range from 0 to 100,
with a mean of 33.24.

The distribution is slightly right-skewed, with a noticeable concentration
of tracks at lower popularity levels. A significant number of tracks
have very low or zero popularity, indicating that only a small portion
of songs achieve high engagement on the platform. This spike at zero is likely attributable to inactive, historical, or 
regional tracks that have accumulated little to no recent streaming 
activity on the platform.

\begin{figure}[H]
\centering
\includegraphics[width=0.8\textwidth]{popularity_dist.png}
\caption{Distribution of song popularity}
\label{fig:pop_hist}
\end{figure}

% -------------------------------------------------------------

\subsection{Danceability vs Popularity}

A scatter plot of danceability against popularity, along with a fitted
linear trend line, was used to examine their relationship.

A weak positive relationship is observed, suggesting that songs with
higher danceability tend to have slightly higher popularity. However,
the large spread of data points indicates that danceability alone is
not a strong predictor of popularity.

\begin{figure}[H]
\centering
\includegraphics[width=0.8\textwidth]{dancebility_pop.png}
\caption{Danceability vs popularity}
\label{fig:dance_pop}
\end{figure}

% -------------------------------------------------------------

\subsection{Energy vs Popularity}

The relationship between energy and popularity was analysed using a
scatter plot with a linear trend line.

The plot shows almost no visible linear trend, which is consistent with
the near-zero correlation value. Although energy does not appear strongly
related to popularity individually, it contributes to the model when
combined with other features. This is a classic example of a suppression effect, where the true 
influence of a variable is masked at the bivariate level but becomes 
statistically significant in the presence of correlated predictors 
such as \texttt{loudness} ($r = 0.762$).

\begin{figure}[H]
\centering
\includegraphics[width=0.8\textwidth]{energy_pop.png}
\caption{Energy vs popularity}
\label{fig:energy_pop}
\end{figure}

% -------------------------------------------------------------

\subsection{Correlation Heatmap}

A correlation heatmap was used to visualise relationships among all key
variables.

Strong positive correlation is observed between \texttt{energy} and 
\texttt{loudness} ($r = 0.762$), while a strong negative correlation 
exists between \texttt{energy} and \texttt{acousticness} ($r = -0.734$). 
The strongest correlation with popularity is observed for 
\texttt{instrumentalness} ($r = -0.095$), which remains notably weak, 
confirming that no single audio feature is a strong linear predictor 
of popularity.

\begin{figure}[H]
\centering
\includegraphics[width=0.8\textwidth]{correlation_heatmap.png}
\caption{Correlation heatmap of audio features}
\label{fig:heatmap}
\end{figure}

% -------------------------------------------------------------

\subsection{Popular vs Non-Popular Songs}

Box plots were used to compare audio features between popular
(popularity $\geq 60$) and non-popular songs (popularity $< 60$).

Popular songs generally show higher danceability and lower
instrumentalness and speechiness. These differences support the
regression findings and suggest that songs with clearer vocals and
better rhythm tend to perform better.

\begin{figure}[H]
\centering
\includegraphics[width=0.8\textwidth]{pop_non_pop.png}
\caption{Comparison of audio features between popular and non-popular songs}
\label{fig:boxplot}
\end{figure}

% -------------------------------------------------------------

\subsection{Model Comparison Visualization}

The comparison of $R^2$ and Adjusted $R^2$ across the three regression
models is shown below.

A clear improvement is observed from Model 1 to Model 2, while Model 3
shows only a slight increase. This indicates that most of the explanatory
power is contributed by the first two blocks of features.

\begin{figure}[H]
\centering
\includegraphics[width=0.8\textwidth]{r2_comparisons.png}
\caption{$R^2$ and Adjusted $R^2$ across regression models}
\label{fig:r2}
\end{figure}

\subsection{Combined Research Dashboard}

All seven visualisations were assembled into a single unified dashboard
using the \texttt{patchwork} package, providing a comprehensive overview
of the analysis. The dashboard presents the popularity distribution,
scatter plots, box plots, model comparison, regression coefficients,
and correlation heatmap in a single layout, making it easier to
interpret the findings holistically.

\begin{figure}[H]
\centering
\includegraphics[width=\textwidth]{spotify_dashboard.png}
\caption{Combined research dashboard — Spotify audio features analysis}
\label{fig:dashboard}
\end{figure}










% =============================================================
\section{Results}
% =============================================================

\subsection{Model Comparison}

Table 5 presents the fit statistics for all three
hierarchical regression models.

\begin{table}[H]
\centering
\small
\renewcommand{\arraystretch}{1.4}
\begin{tabular}{@{}lcccccc@{}}
\toprule
\textbf{Model} 
  & \textbf{$R^2$} 
  & \textbf{Adj.\ $R^2$} 
  & \textbf{AIC} 
  & \textbf{$F$} 
  & \textbf{df} 
  & \textbf{$p$} \\
\midrule
Model 1 (Block 1) & 0.0057 & 0.0056 & 1{,}030{,}778 & 216 & (3, 113996)   & $<0.001$ \\
Model 2 (Block 1+2) & 0.0218 & 0.0218 & 1{,}028{,}917 & 424 & (6, 113993) & $<0.001$ \\
Model 3 (Full)      & 0.0224 & 0.0223 & 1{,}028{,}857 & 326 & (8, 113991) & $<0.001$ \\
\bottomrule
\end{tabular}
\vspace{0.3cm}

{\centering \textit{Table 5: Hierarchical regression model fit statistics.
Block 1 = danceability, energy, valence.
Block 2 = acousticness, instrumentalness, speechiness.
Full model adds tempo and loudness.}\par}
\label{tab:models}
\end{table}

All three models were statistically significant, indicating that the
selected audio features collectively contribute to explaining variations
in song popularity.

The addition of Block 2 (acoustic texture features) led to a substantial
increase in explanatory power. The $R^2$ value increased from 0.0057 to
0.0218, representing an improvement of 0.0161. This increase was highly
significant, as confirmed by the ANOVA test
($F(3, 113993) = 627.87$, $p < 2.2 \times 10^{-16}$).

The addition of Block 3 (tempo and loudness) resulted in only a marginal
increase in $R^2$ to 0.0224. Although this improvement was statistically
significant ($F(2, 113991) = 32.03$, $p = 1.24 \times 10^{-14}$),
its practical contribution was relatively small.

Overall, these results indicate that most of the explanatory power comes
from the first two blocks of features, while additional variables provide
diminishing returns.

% -------------------------------------------------------------

\subsection{Coefficient Analysis --- Full Model}

Table 6 presents the regression coefficients for the full
model (Model 3), including estimates, standard errors, $t$-values,
$p$-values, and 95\% confidence intervals.

All predictors in the model were found to be statistically significant.

\begin{table}[ht]
\centering

\begin{tabular}{@{}lrrrrrr@{}}
\toprule
\textbf{Predictor} & \textbf{Est.} & \textbf{SE}
  & \textbf{$t$} & \textbf{$p$}
  & \textbf{95\% CI lower} & \textbf{95\% CI upper} \\
\midrule
(Intercept)       &  36.10 & 0.656  &  54.9  & $< 0.001$ &  34.80 &  37.30 \\
danceability      &   8.91 & 0.455  &  19.6  & $< 0.001$ &   8.02 &   9.80 \\
energy            &  $-$2.33 & 0.524 & $-$4.44 & $< 0.001$ & $-$3.36 & $-$1.30 \\
valence           &  $-$9.64 & 0.313 & $-$30.8 & $< 0.001$ & $-$10.30 & $-$9.03 \\
acousticness      &  $-$0.97 & 0.302 & $-$3.20 & 0.0014    & $-$1.56  & $-$0.38 \\
instrumentalness  &  $-$8.31 & 0.254 & $-$32.8 & $< 0.001$ & $-$8.80  & $-$7.81 \\
speechiness       & $-$12.00 & 0.645 & $-$18.6 & $< 0.001$ & $-$13.20 & $-$10.70 \\
tempo             &   0.014 & 0.002 &   6.19 & $< 0.001$ &   0.010  &   0.019 \\
loudness          &   0.110 & 0.024 &   4.69 & $< 0.001$ &   0.064  &   0.156 \\
\bottomrule
\end{tabular}

\vspace{0.3cm}

{\centering \textit{Table 6: OLS regression coefficients for the full model (Model 3), $n = 114{,}000$}\par}

\label{tab:coeff}

\end{table}
\begin{figure}[H]
\centering
\includegraphics[width=0.8\textwidth]{coeffecients.png}
\caption{Regression coefficients with 95\% confidence intervals (Model 3)}
\label{fig:coeff}
\end{figure}

Among the predictors, \textbf{danceability} showed the strongest positive
association with popularity, indicating that more danceable songs tend
to perform better.

In contrast, \textbf{speechiness}, \textbf{valence}, and
\textbf{instrumentalness} exhibited strong negative relationships with
popularity. This suggests that songs with excessive spoken content,
higher emotional positivity, or fewer vocals tend to perform less well.

\textbf{Loudness} and \textbf{tempo} had small but positive effects,
indicating that slightly louder and faster songs may have a marginal
advantage.

Despite all predictors being statistically significant, the overall
$R^2$ value (0.0224) remains low. This indicates that audio features
alone explain only a small portion of the variation in song popularity,
suggesting the influence of external factors such as marketing, artist
popularity, and release timing.
% =============================================================
\section{Inferences and Insights}
% =============================================================

The results of the analysis provide several important insights into how
different audio features relate to song popularity.

\begin{itemize}

  \item \textbf{Danceability is the strongest positive predictor.}
        A one-unit increase in danceability (on a 0--1 scale) is associated
        with an increase of approximately 8.91 points in popularity
        ($p < 0.001$, 95\% CI [8.02, 9.80]). This indicates that
        rhythm-focused and dance-friendly tracks tend to attract more
        listeners and perform better on streaming platforms.

  \item \textbf{Speechiness has the strongest negative impact.}
        Speechiness shows the largest negative effect
        ($\hat{\beta} = -12.00$, 95\% CI $[-13.2, -10.7]$), suggesting
        that tracks with higher spoken-word content are generally less
        popular. This may be because the popularity metric mainly reflects
        music consumption rather than spoken content.

  \item \textbf{Instrumental tracks are less popular.}
        Instrumentalness has a strong negative association with popularity
        ($\hat{\beta} = -8.31$, $p < 0.001$), indicating that songs without
        vocals tend to receive fewer streams compared to vocal-based tracks.
        This reflects the dominance of vocal music in mainstream listening.

  \item \textbf{Valence shows an unexpected negative relationship.}
        Valence, which represents musical positivity, has a negative effect
        ($\hat{\beta} = -9.64$, $p < 0.001$). This suggests that highly
        positive or cheerful songs are not necessarily more popular.
        Emotionally complex or slightly melancholic tracks may engage
        listeners more deeply, leading to higher popularity.

  \item \textbf{Energy and acousticness have smaller but significant effects.}
        Energy shows a moderate negative effect ($\hat{\beta} = -2.33$,
        $p < 0.001$), while acousticness has a small negative effect
        ($\hat{\beta} = -0.97$, $p = 0.0014$). These features contribute
        to popularity but are not strong individual predictors.

  \item \textbf{Tempo and loudness show positive contributions.}
        Both tempo ($\hat{\beta} = 0.014$, $p < 0.001$) and loudness
        ($\hat{\beta} = 0.110$, $p < 0.001$) have positive coefficients,
        indicating that faster and slightly louder tracks tend to perform
        marginally better. However, their impact is relatively small
        compared to other features.

  \item \textbf{Overall explanatory power is limited.}
        The overall $R^2$ value of 0.0224 indicates that the model explains
        only 2.24\% of the variation in popularity. This highlights that
        audio features alone are not sufficient to determine a song's
        success. Other factors such as artist popularity, marketing
        strategies, playlist placement, and social media influence likely
        play a much more significant role in driving streaming performance
        \citep{interiano2018}.

\end{itemize}

% =============================================================
\section{Conclusion}
% =============================================================

This study applied a hierarchical OLS regression framework in R to
analyze the relationship between Spotify audio features and song
popularity across a dataset of 114{,}000 tracks.

The results show that all selected audio features are statistically
significant predictors of popularity. The inclusion of additional
feature blocks led to measurable improvements in model performance,
as confirmed by highly significant ANOVA tests
($F(3, 113993) = 627.87$, $p < 2.2 \times 10^{-16}$ for Block 2;
$F(2, 113991) = 32.03$, $p = 1.24 \times 10^{-14}$ for Block 3). However, despite
this statistical significance, the overall explanatory power of the
model remains limited, with an $R^2$ value of 0.0224.

Among the features, danceability emerged as the strongest positive
predictor, while speechiness, instrumentalness, and valence showed
strong negative relationships with popularity. These findings suggest
that songs with strong rhythm, clear vocals, and balanced emotional
tones are more likely to perform well on streaming platforms.

At the same time, the relatively low $R^2$ highlights an important
insight: audio features alone cannot fully explain song popularity.
External factors such as artist reputation, marketing strategies,
playlist placement, and audience trends likely play a much more
significant role \citep{interiano2018}.

Future work incorporating such variables alongside audio features 
could substantially improve predictive accuracy and provide a 
more complete picture of what drives success in the streaming era.

Overall, this study demonstrates the usefulness of data-driven
approaches in understanding music trends, while also emphasizing
the complexity of predicting popularity in the digital music
ecosystem.

% =============================================================
\section{Future Scope}
\label{sec:future}
% =============================================================

The present study can be extended in several meaningful directions to
improve both the depth and accuracy of the analysis.

\begin{enumerate}

  \item \textbf{Genre-wise Analysis.}  
  Performing separate analyses for different music genres (such as pop,
  rock, hip-hop, and classical) can help determine whether the influence
  of audio features on popularity varies across genres. The \texttt{track\_genre} variable available in the dataset 
provides a natural starting point for such an analysis.

  \item \textbf{Advanced Machine Learning Models.}  
  Using non-linear models such as decision trees, random forests, or
  gradient boosting techniques can help capture complex relationships
  and interactions between variables that are not detected by linear
  regression.

  \item \textbf{Temporal Trend Analysis.}  
  Incorporating time-based data can help analyze how music preferences
  and feature importance have evolved over the years. This would provide
  insights into changing listener behavior and industry trends.

  \item \textbf{Inclusion of External Factors.}  
  Adding variables such as artist popularity, playlist inclusion,
  release timing, and marketing influence can significantly improve
  the predictive power of the model and provide a more realistic
  understanding of song success. Given that the full model achieved an $R^2$ of only 0.0224, 
the inclusion of such external variables could substantially 
improve explanatory power.

  \item \textbf{Feature Engineering and Dimensionality Reduction.}  
  Techniques such as Principal Component Analysis (PCA) can be used to
  reduce redundancy among correlated variables (e.g., energy and
  loudness) and identify underlying patterns in the data.
\item \textbf{Integration with User Behaviour Data.}
        Combining audio features with user interaction data such as
        skip rates, playlist saves, and listening duration could
        provide a richer understanding of what drives engagement
        beyond popularity scores alone.
\end{enumerate}

% =============================================================
\section{R Implementation (Code Snippets)}
% =============================================================

This section presents key code snippets used in the analysis. The complete
R script is provided as supplementary material.

\subsection{Data Loading and Cleaning}

\begin{lstlisting}
library(tidyverse)
library(janitor)

df <- read_csv("spotify_tracks.csv") |>
      clean_names() |>
      drop_na(popularity, danceability, energy, valence,
              acousticness, instrumentalness, speechiness,
              tempo, loudness) |>
      distinct()
\end{lstlisting}   


\subsection{Hierarchical Regression Models}

\begin{lstlisting}  
model1 <- lm(popularity ~ danceability + energy + valence,
             data = df)

model2 <- lm(popularity ~ danceability + energy + valence +
               acousticness + instrumentalness + speechiness,
             data = df)

model3 <- lm(popularity ~ danceability + energy + valence +
               acousticness + instrumentalness + speechiness +
               tempo + loudness,
             data = df)

anova(model1, model2, model3)
\end{lstlisting}


\subsection{Model Comparison Metrics}

\begin{lstlisting}
extract_fit <- function(model, label) {
  s <- summary(model)
  tibble(
    Model  = label,
    R2     = s$r.squared,
    Adj_R2 = s$adj.r.squared,
    AIC    = AIC(model)
  )
}

model_comparison <- bind_rows(
  extract_fit(model1, "Model 1"),
  extract_fit(model2, "Model 2"),
  extract_fit(model3, "Model 3")
)
\end{lstlisting}


\subsection{ANOVA Model Comparison}

\begin{lstlisting}
anova_result <- anova(model1, model2, model3)
print(anova_result)
\end{lstlisting}


\subsection{Coefficient Extraction}

\begin{lstlisting}
library(broom)

coeff_table <- tidy(model3,
                    conf.int = TRUE,
                    conf.level = 0.95)
\end{lstlisting}


\subsection{Popular vs Non-Popular Classification}

\begin{lstlisting}
df <- df |>
  mutate(popular_cat = ifelse(popularity >= 60,
                              "Popular",
                              "Non-Popular"))
\end{lstlisting}


\subsection{Data Visualization (Example)}

\begin{lstlisting}
ggplot(df, aes(x = danceability, y = popularity)) +
  geom_point(alpha = 0.2, color = "#1DB954") +
  geom_smooth(method = "lm", color = "black") +
  labs(title = "Danceability vs Popularity")
\end{lstlisting}

\subsection{Correlation Heatmap}

\begin{lstlisting}
library(corrplot)

corr_matrix <- cor(df |> select(popularity, danceability, energy,
                               valence, acousticness,
                               instrumentalness, speechiness,
                               tempo, loudness))

corrplot(corr_matrix, method = "color", type = "upper")
\end{lstlisting}


\subsection{Combined Dashboard}

\begin{lstlisting}
library(patchwork)
dashboard <- (p1 | p2 | p3) /
             (p4 | p5) /
             (p6 | p7)
ggsave("spotify_dashboard.png", dashboard,
        width = 18, height = 16, dpi = 150)
\end{lstlisting}


\subsection{Saving Output Files}

\begin{lstlisting}
write_csv(df, "spotify_clean.csv")
write_csv(model_comparison, "model_comparison.csv")
write_csv(coeff_table, "coefficients.csv")
\end{lstlisting}

% =============================================================
% =============================================================

% =============================================================

\bibliographystyle{plainnat}
\bibliography{RJreferences}
% =============================================================
%  AUTHOR ADDRESS
% =============================================================

\end{document}
